% Supplementary material: knowledge-distillation detail moved from Section 4.2 of the main text.

\section{Knowledge Distillation Variants}
\label{supp:kd}

\subsection{White-Box vs.\ Black-Box KD}

The survey by \citet{yang2024kd_survey} classifies \gls{llm} distillation into two paradigms:

\emph{White-box \gls{kd}} has access to the teacher's intermediate representations. TinyBERT \citep{jiao2020tinybert} pioneered layer-by-layer feature alignment, matching student hidden states, attention matrices, and embedding layers to teacher counterparts via MSE loss. In the decoder-only \gls{llm} setting, this requires pairing student and teacher layers, which is non-trivial when sizes differ substantially, and is computationally expensive but yields superior retention of teacher capabilities.

\emph{Black-box \gls{kd}} uses only the teacher's generated outputs, either token probabilities or free-form text, as the supervision signal. Orca \citep{mukherjee2023orca} exemplifies text-output black-box \gls{kd}: GPT-4 generates explanation traces and step-by-step reasoning, which a 13B LLaMA-based student learns to imitate. DeepSeek-R1 \citep{deepseek2025r1} extends this to long chain-of-thought reasoning distillation, fine-tuning 1.5B--70B Qwen and Llama students on the reasoning traces of its large model trained by reinforcement learning; its authors report that the reasoning patterns of the large model can be transferred to smaller ones in this way (\evC).

\subsection{Chain-of-Thought Distillation}

A particularly productive form of black-box \gls{kd} is \gls{cot} distillation \citep{wei2022cot}: the teacher model generates explicit reasoning steps, and the student learns to produce or internalize these steps. \citet{feng2024counterfactual} demonstrate counterfactual distillation, which augments \gls{cot} traces with contrastive incorrect reasoning, as a way to sharpen the student's reasoning boundaries. \citet{liao2024nesycd} propose Neural-Symbolic Collaborative Distillation, using a neuro-symbolic teacher to provide structured reasoning supervision that is more transferable to small students than purely neural \gls{cot}.

\subsection{Data Augmentation in KD}

The survey by \citet{xu2024kd_survey} highlights data augmentation (DA) as a central paradigm within \gls{kd} frameworks. When the student lacks sufficient coverage of the teacher's knowledge space, the teacher can generate synthetic training instances targeting specific skill gaps, the mechanism underlying the Phi family's success \citep{gunasekar2023phi1}. \citet{ballout2024efficient_kd} use a roughly 3B teacher's most influential input tokens, identified by attribution methods such as saliency maps, as rationales for the student, and report improvements over standard fine-tuning and over prior distillation methods on four datasets.

\subsection{Self-Distillation}

An emerging approach eliminates the external teacher entirely. \citet{fu2024dynamic_selfd} propose Dynamic Self-Distillation, in which the model uses its own predictions from previous mini-batches as soft targets during fine-tuning, providing a form of implicit regularization that reduces overfitting on small fine-tuning datasets, a critical challenge for \glspl{slm}.

% ----------------------------------------------------------------
